\HMTNarrativeHeading{La forme conservée par la seconde projection}

La réalisation exceptionnelle développe l’incidence qui accompagne
la génération numérique. Chaque préfixe transmet des relations de frontière
compatibles avec les troncatures. L’encodage de Paley conserve
le mot et son image ; le plan de Witt organise leurs supports ; le
relèvement vers Golay introduit l’incidence binaire avec ses marques.
Le parcours ultérieur construit les réalisations réticulaires et
opératorielles correspondantes.

Les étoiles d'incidence permettent d'observer ce lien. Le registre de couplage détermine un quadruplet. Les hexades qui le contiennent distinguent quatre paires complémentaires et une signature sélectionne l'hexade marquée. Dans la figure, les réunions représentent des opérations entre ensembles. L'élévation aux octades conserve cette configuration et distingue une cinquième octade transversale. Les cinq régions numériques maintiennent ainsi une correspondance avec des relations d'incidence spécifiées.

L’égalité d’une lecture ternaire peut réunir des régions dont
l’orientation, la multiplicité et la position restent distinctes. Ces
données interviennent dans la correspondance exceptionnelle. La seconde
projection conserve précisément des relations que la lecture de valeur
traite autrement. La compatibilité entre les deux s’exprime sur
l’état enrichi commun et ses applications de raffinement.

La théorie M et la dualité T sont présentées au moyen des réalisations
dimensionnelles et de leurs conditions de compatibilité. L’action du
Monstre a son porteur dans l’espace construit pour elle. Les
lecteurs numériques, les configurations finies, les réseaux et les
représentations reçoivent des domaines distincts au sein d’une même
généalogie. Cette précision permet de suivre comment une structure
incidencielle se transmet jusqu’à ses conséquences ultérieures.

La partie suivante conserve les constructions, les tableaux et les preuves
qui rendent le parcours effectif. La narration ramène chaque
réalisation à la question directrice : quelles relations survivent à la
projection et comment elles se récupèrent dans l’objet auquel elles parviennent. La
conservation d’une forme de Gram, la covariance d’un encodage
et l’invariance d’une action auront leurs preuves propres. Leur
articulation apporte le contenu géométrique de la conservation
évolutive de l’unité.

